\section{The global inequality on an inert all-row declaration}

% BEGIN G82I1
\paragraph{G82I1. An inert index in terms of its row lengths.}
Fix a late member with the hub of F78D2, a twisted middle-range
index $i$ outside the type primes, $p=\ell_i$ inert, and
$\pi\in Z_{\rm good}(i)$. Put $v=D.\mathrm{nd.vdel}(i)$,
$c_B=X.c_B(i)$, and $d_w=[E_w:\mathbb Q_p]$. Suppose for each
$w\mid p$ there is a real $B_w$ bounding the normalized length
of the finite integral cokernel for all the dual row maps of
I81G3 there. A bound on the entire native row test is sufficient.
Then one has
\[
\begin{aligned}
 v_{\rm CC}(i,\pi)\le {}&
     s_\pi+\sum_{0<j<t}u_j+v_p(h_E)+p^{-2}\\
 &+\sum_{w\mid p} B_w+
       t c_B\sum_{w\mid p}(d_w+1)+2(n_V-1)v .
\end{aligned}
\]
Here $u_j$ is the native cyclotomic length for the ratio
$\chi_{j+1}/\chi_1$. We verify the nonunit discriminant term.

By I82W1 the completions at $p$ are unramified, $p\ge11$, and the
native level exponent is one because $p$ is outside the small
conductor. Local inertia contains an element with cyclotomic
value $1+p$ by total ramification of cyclotomic layers over an
unramified field. Its character values at $\pi$ are distinct,
and their root-difference product at the first character has
valuation $\sum_{1\le j<t}(1+v_p(j))=v$. Normalize the matrices
at $\pi$ and perturb the character corner at the other points
as in G81S3. Its hypotheses on the Hecke polynomials, full
rank matrix block and off-support fixed vectors are supplied by
D79G1--D79G3 and D79G7 at this good localization. The argument
excluding repeated global constituents uses the simple labelled
jumps of the full family in S80J1 and crystallinity from S80J3.
It applies when the successive letters have different CM types.
After a good finite coefficient enlargement the ordered-root basis,
integral first-character matrix $E$ and nonzero differences at
all points are therefore available, including at primes $p\le t$.

The inert declaration has no columns, by $\mathrm{DataAD.col}$.
Use G82C1 with its family of extra generators empty, so its
module is simply $Y=\sum_g A(U_g-1)E$. At each lift off support
the line $E(x')V$ gives a nonzero homomorphism $Y\to V$ exactly
as in G82C2; there are no local relations there to impose.
G82C2 thus bounds the normalized length of the evaluation ideal
quotient by the sum of the normalized lengths of the block
homomorphism spaces. Each block map is determined by its actual
cocycle $g\mapsto\varphi((U_g-1)E)$. G81S3 bounds the total
length lost in passing to global classes by $2(n_V-1)v$ since
the first character has zero coboundaries. This is the only
cost from the root differences in this calculation.

Every rank class has the common graph-lattice description of
I81G3. Simultaneous restriction modulo the finite parts at all
places over $p$ costs at most $\sum_w B_w$ by the local Tate
pairing of I81G4. Finite lists suffice by I81G2. The rank
kernel is in the chosen Selmer group by the away-from-$p$
conditions of D79G7 and G79E8. For each character impose the
auxiliary inertia conditions at $Q_i$ and impose the finite
condition at each $w\mid p$. M81B1--M81B3 supply the local
character test with the stored $c_B$ at these inert places, so
I81G5 bounds its normalized row loss by $(d_w+1)c_B$.
G79E8 bounds the auxiliary loss for each character by
$\sum_{P\in Q_i} z_P/k_O$ on the common coefficient DVR. The
twisted-index killing clause bounds $t$ times this sum by $p^{-2}$.
The remaining character classes are unramified everywhere for
the trivial ratio and locally zero above $p$ for all nontrivial
ratios. G79E6 gives their bounds $v_p(h_E)$ and $u_j$.
The original finiteness guards apply at the good point. I81G2
and G81S3 descend the normalized Selmer, local, auxiliary and
evaluation lengths to the native coefficients. Adding yields
the asserted inequality.
% END G82I1

% BEGIN G82I2
\paragraph{G82I2. The exact allowance and the currently covered inert indices.}
Keep the point of G82I1, and put $m_p=|\{w:w\mid p\}|$.
The cover condition and the empty column set imply that the
row set consists of all the places over $p$. With the prescribed
hub lengths $b_w=X.\mathrm{lenRow}(i,w)$, its two terms in
the target inequality are therefore
\[
\begin{aligned}
 \mathrm{sumOver}(D,p,X.\mathrm{len0}(i,\pi))&=\sum_w b_w,\\
 X.\mathrm{errorBoundHub}(i,\pi)
               &=t c_B\sum_w(d_w+2)+4n_Vv .
\end{aligned}
\]
These equalities are the inert $\mathrm{len0I}$ and
$\mathrm{error\_bound\_inert\_eq}$ fields of $\mathrm{DataAD}$;
the final column-invariant term has coefficient zero. Consequently
the native $\mathrm{StarAD}$ inequality at this point follows
from G82I1 whenever $B_w\le b_w$ at every place. More generally
it follows if
\[
           \sum_w B_w\le\sum_w b_w+t c_B m_p+(2n_V+2)v ,
\]
by subtracting the character and global-coboundary terms of
G82I1 from the two displayed allowances. This last comparison
allows savings proved at one place to compensate for an excess
at another. The choice $B_w=b_w$ suffices throughout.

In row B the inequality thus holds at every guarded index where
all the orbits satisfy I82B1's numerical condition
$25h(w)/p+3\le3C_2 U$ with its stated orbit height. This includes
its completely nonthick local case. Outside row B it holds at
every guarded index in the prime range of N82I2, even in the
presence of thick letters. N82I1 also supplies the requisite
length at every completely nonthick orbit outside that prime
range. In particular for an index outside row B whose every
orbit is covered by either of those tests, the global inequality
holds for every good point. These assertions follow by applying
the proved local bounds before imposing the simultaneous global
conditions. R82W4's lifts of Wach vectors at cuts by themselves
do not furnish the remaining integral cohomology row lengths.

Finally the host classification on $\mathrm{RangeAD}$ outside
the type primes gives split or inert, since the small primes and
$q,d$ are excluded. The split branch at every guarded point is
G82C5. Thus supplying the above total comparison at the remaining
inert points would give the full $\mathrm{StarAD}$ field
$\mathrm{Concl.mid}$ on the same hub. The numerical size of
$X.\mathrm{errorBoundHub}$ already has its separate family bound
in J81H2. The following chapters supply the row B lengths at all guarded
points for a common choice of row constant and further local
and aggregate comparisons in rows C and D.
% END G82I2
